\section{Related work}
\label{sec:related_work}

\subsection{Image complexity as a perceptual and computational target}

Visual complexity has been operationalised through element count, organisation, search or encoding effort, edges, entropy, compression, clutter, and learned semantic representations \citep{donderi2006visual,rosenholtz2007measuring,forsythe2011predicting}. Web-page, intermediate-feature, and attention studies further show that both domain and task alter what a complexity score captures \citep{wu2013webcomplexity,saraee2020complexity,sun2021curious}. These alternatives are not interchangeable. IC9600 combines 9,600 images, multi-observer scores, and local complexity maps, and demonstrates the value of both detail and contextual branches \citep{feng2023ic9600}. Kyle-Davidson et al. show that low-level properties and semantic information both contribute to perceived scene complexity \citep{kyledavidson2023complexity}. Segment- and class-count models similarly emphasise the importance of interpretable semantic structure across datasets \citep{shen2024simplicity}. Recent JVCIR work explicitly aligns global complexity prediction with regional localisation and eye-tracking evidence \citep{guo2026dual}.

Large-scale and personalised aesthetics datasets such as AVA, AADB, PARA, TAD66K, and LAPIS likewise demonstrate that attribute names and predictive claims require explicit annotation scope \citep{murray2012ava,kong2016photo,yang2022para,he2022tad66k,maerten2025lapis}. Recent colour-aesthetics benchmarking and self-supervised or meta-learning work further separate dataset, task, and representation choices \citep{he2023color,yang2024selfsupervised,yan2024hybrid}. General-purpose learned representations such as CLIP and DINOv2 can supply rich image embeddings \citep{radford2021clip,oquab2024dinov2}, but their reproducible extraction still does not by itself define the target construct. Multiscale structural complexity evaluated on controlled stimulus sets \citep{kravchenko2026multiscale} and pre-registered studies of visualization-image complexity with multiple interpretable metrics \citep{chu2026visualization} further demonstrate that a single composite score rarely captures all intended variance, reinforcing the need for transparent component-level diagnostics of the kind we provide.

Those studies reinforce the boundary of the present work. A human-perception claim requires human evidence and natural-image diversity. Our no-human synthetic targets instead provide controlled interventions for debugging image-analysis software. The named legacy formulas are subjects of a failure analysis, not competitors to IC9600 or recent learned methods.

\subsection{Handcrafted descriptors and saliency dependencies}

Classical descriptors remain useful because their operations can be inspected: Canny edge fraction, image entropy, compressed size, mirror similarity, colourfulness \citep{hasler2003colourfulness}, contrast statistics \citep{peli1990contrast}, and explicit layout attributes \citep{ngo2003interface}. Explicit algebra, however, does not guarantee an active implementation. A formula can list a saliency feature while the runtime supplies a constant.

Saliency itself spans distinct tasks and algorithms. Itti et al. model rapid visual attention from multiscale feature maps \citep{itti1998saliency}. Hou and Zhang derive a fast saliency map from the residual between an image log spectrum and its local spectral average \citep{hou2007spectral}. We use the latter as a deterministic, dependency-minimal repair because it can be implemented with NumPy FFT and OpenCV core. It is an operational spectral-residual map, not an eye-fixation or salient-object ground truth.

\subsection{Reproducibility, semantic health, and measurement}

Reproducibility programmes rightly emphasise fixed code, environments, seeds, complete artefacts, and strong baselines \citep{pineau2021reproducibility}. These controls answer whether the same computation can be reconstructed. Semantic health asks a complementary question: did the computation retain the variation and distinctions required by its declared role? A constant weighted feature, a score dominated by clipping, or two sign-reversed baseline candidates can all be perfectly reproducible.

This distinction parallels construct-validity guidance. A name and a favourable correlation do not establish a measure; evidence must connect operations, constructs, alternatives, and observable consequences \citep{cronbach1955construct,flake2020measurement}. Our gates are narrower than construct validation: they can reject an inactive computation, but cannot prove that an active computation measures a human construct. This asymmetry is deliberate and makes the checks suitable for continuous integration.

\subsection{Testing scientific and image-processing software}

The oracle problem---the absence of a known-correct output for a given input---is central to testing scientific software \citep{guderlei2007statistical,ding2010selfchecked}. Metamorphic testing addresses this by defining relations between outputs of related inputs: if an image is rotated, an edge detector's total edge count should not change by more than a bounded factor \citep{lin2018metamorphic}. Property-based testing generalises this idea to arbitrary input generators and declared invariants. Recent work specifically targets image-processing applications: MT4Image combines metamorphic relations with coverage analysis for deep-learning image pipelines \citep{mt4image2025}.

Our semantic-health gates share the oracle-free philosophy of metamorphic testing but target a different failure class. Metamorphic relations typically detect functional errors visible in output transformations. The failure studied here---a constant feature column that passes all finite-value, row-count, and hash checks---is invisible to output-level metamorphic relations unless the relation explicitly requires non-degenerate variation in intermediate representations. The responsiveness gate (every nonzero-weight feature must vary across diverse non-uniform images) and the retained-weight gate (the formula's effective dimensionality must match its declared dimensionality) operate at the intermediate-representation layer rather than the input--output layer. We position these gates as complementary to, not a replacement for, metamorphic and property-based testing of visual pipelines.

\subsection{Synthetic stress tests and renderer shift}

Synthetic data permit controlled factors and exact manifests, but shared primitives and renderer assumptions limit external validity \citep{tobin2017domainrandomization,torralba2011datasetbias}. We therefore separate three claims: recovery within the Pillow generator, response under a separately implemented internal OpenCV renderer, and generalisation to natural images. Only the first two are examined. Complete cross-target tables and mechanism descriptors are used to reveal confounding rather than to label the second renderer ``external validation.''
